\documentclass{article}%
\usepackage[T1]{fontenc}%
\usepackage[utf8]{inputenc}%
\usepackage{lmodern}%
\usepackage{textcomp}%
\usepackage{lastpage}%
\usepackage{geometry}%
\geometry{margin=2cm}%
\usepackage{expex}%
\usepackage[T1]{fontenc}%
\usepackage[utf8]{inputenc}%
\usepackage[spanish]{babel}%
%
\title{Análisis Lingüístico: Gramatica Normativa Kaqchikel}%
\author{Generado por el TFM de Elías Carrasco}%
%
\begin{document}%
\normalsize%
\maketitle%
\section{Page 51}%
\label{sec:Page51}%
a) %
\subsection{Sintagma Nominal}%
\label{subsec:SintagmaNominal}%

%
En el idioma K'iche' se presentan los siguientes casos de sintagma nominal, con los %
siguientes modificadores: demostrativos, adjetivos, clasificadores personales, pronombres, %
afectivos, pluralizadores, negativos y poseedores. %
\subsection{Demostrativo}%
\label{subsec:Demostrativo}%

%
Se refieren a un nominal identificado o conocido. %
\par\medskip%
\ex
\textit{Vendí el caballo} \\
\begingl
\gla Xink'ayij le kej. //
\glb Xink'ayij le kej //
\glc \textbf{verb} \textbf{DEM} \textbf{caballo} //
\glft \textbf{Sintaxis:} [le kej]\textsubscript{SN} //
\endgl
\xe
%
\medskip%
\par\medskip%
\ex
\textit{Regalaron la faja} \\
\begingl
\gla Xkisipaj we pas. //
\glb Xkisipaj we pas //
\glc \textbf{verb} \textbf{DEM} \textbf{faja} //
\glft \textbf{Sintaxis:} [we pas]\textsubscript{SN} //
\endgl
\xe
%
\medskip%
\par\medskip%
\ex
\textit{Trabajamos la tierra.} \\
\begingl
\gla Kqachakuj ri ulew. //
\glb Kqachakuj ri ulew //
\glc \textbf{verb} \textbf{DEM} \textbf{tierra} //
\glft \textbf{Sintaxis:} [ri ulew]\textsubscript{SN} //
\endgl
\xe
%
\medskip%
\subsection{Adjetivo}%
\label{subsec:Adjetivo}%

%
El adjetivo va en posición anterior a los núcleos nominales. %
\par\medskip%
\ex
\textit{vean un güipil bonito.} \\
\begingl
\gla Chiwila pe' jun je'l po't. //
\glb Chiwila pe' jun je'l po't //
\glc \textbf{verb} \textbf{particle} \textbf{ART} \textbf{bonito} \textbf{güipil} //
\glft \textbf{Sintaxis:} [jun je'l po't]\textsubscript{SN} //
\endgl
\xe
%
\medskip%
\par\medskip%
\ex
\textit{¡Que hermosa se ve esa ropa roja!} \\
\begingl
\gla Je'l kwil le kaq atz'yaq. //
\glb Je'l kwil le kaq atz'yaq //
\glc \textbf{hermosa} \textbf{verb} \textbf{DEM} \textbf{roja} \textbf{ropa} //
\glft \textbf{Sintaxis:} [le kaq atz'yaq]\textsubscript{SN} //
\endgl
\xe
%
\medskip%
\subsection{Clasificadores}%
\label{subsec:Clasificadores}%

%
\par\medskip%
\ex
\textit{La señorita Xpuch se está riendo.} \\
\begingl
\gla Ktze'n le ali Xpuch. //
\glb Ktze'n le ali Xpuch //
\glc \textbf{verb} \textbf{DEM} \textbf{señorita} \textbf{name} //
\glft \textbf{Sintaxis:} [le ali Xpuch]\textsubscript{SN} //
\endgl
\xe
%
\medskip%
\par\medskip%
\ex
\textit{Don Juan es una persona muy buena.} \\
\begingl
\gla Ri tat Xwan utzalaj winaq. //
\glb Ri tat Xwan utzalaj winaq //
\glc \textbf{DEM} \textbf{Don} \textbf{name} \textbf{buena} \textbf{persona} //
\glft \textbf{Sintaxis:} [Ri tat Xwan]\textsubscript{SN} //
\endgl
\xe
%
\medskip%
\par\medskip%
\ex
\textit{Juan es un muchacho muy trabajador.} \\
\begingl
\gla A Xwan sib'alaj ajchak. //
\glb A Xwan sib'alaj ajchak //
\glc \textbf{classifier} \textbf{name} \textbf{muy} \textbf{trabajador} //
\glft \textbf{Sintaxis:} [A Xwan]\textsubscript{SN} //
\endgl
\xe
%
\medskip

%
\end{document}